\section{Conclusions}

We have assembled paired single-nucleus and spatial transcriptomic profiles across the full
lung adenocarcinoma invasion sequence and used them to ask what changes as invasion
proceeds, and whether that change can be pushed back \emph{in silico}. Four conclusions
survive the checks we applied.

\begin{enumerate}\itemsep2pt
\item The tissue resolves into seven recurring niche domains, six of them ordinary lung
structures; only one is confined to invasive disease, while lymphoid aggregates are present
in precursor stages.

\item The ECM/interstitial programme that distinguishes the invasive domain is elevated in
invasive disease within patients (20 of 22 paired cases), and the result does not depend on
the genes we added by hand.

\item The spatial copy-number signal lies in the alveolar type II compartment and is
present before invasion, consistent with AT2 as the cell of origin.

\item A query signature derived from a patient-internal disease axis and filtered on
independent survival evidence returns a reproducible set of 34 candidate compounds,
dominated by PI3K/mTOR inhibitors, acting on the ECM/interstitial and AT2 compartments
rather than on the tissue as a whole.
\end{enumerate}

Two features of the work constrain how these should be read. First, sequencing depth and
tissue density are the same measurement in this cohort, which is why several initially
positive results did not survive depth matching and why the remaining claims are stated
narrowly. Second, the candidate compounds are computational hypotheses: they act in cancer
cell lines, we have not tested them in lung fibroblasts or in primary tissue, and their
selection rests on a query whose two libraries share a common structure and therefore
provide robustness rather than independent corroboration. The analysis we would prioritise
next is a within-section comparison of the invasive front against distant regions of the
same section, which holds depth fixed by construction and would remove the dominant
confound from the design.
